\section{Discussion}
\label{sec:discussion}

\subsection{Implications of Planetary Magnetic Dichotomy for Martian Habitats}
The fundamental insight revealed by our global analysis is that future human missions to Mars will face two radically divergent magnetic regimes depending on landing site selection:

\begin{enumerate}
    \item \textbf{The Northern Plains / Low-Field Paradigm ($< \SI{20}{\nano\tesla}$):} Primary candidate sites for initial crewed landings (e.g., Arcadia Planitia and Deuteronilus Mensae) are chosen based on resource availability—specifically shallow, accessible water ice for In Situ Resource Utilization (ISRU) and propellant production. However, our findings demonstrate that these sites represent extreme near-null hypomagnetic environments ($>2,500\times$ to $>5,000\times$ weaker than Earth's GMF). Bioregenerative life support systems (BLSS), crop cultivation chambers, and astronauts residing in northern habitats will operate under chronic, unmitigated magnetic deprivation.
    \item \textbf{The Southern Highlands / Mini-Magnetosphere Outposts ($> \SI{1500}{\nano\tesla}$):} In contrast, scientific outposts sited in the Terra Sirenum or Terra Cimmeria anomaly belts reside within intense, localized crustal magnetic fields. These fields form mini-magnetospheres extending hundreds of kilometers into the ionosphere \citep{Mitchell2001}, capable of partially deflecting low-energy solar wind ions and creating localized plasma cusps. While these fields remain $\sim 25\times$ to $50\times$ below Earth's GMF, they provide a significantly elevated baseline compared to the northern plains.
\end{enumerate}

\subsection{Synergistic Hazards: Hypomagnetism, Hypogravity, and Radiation}
On the Martian surface, hypomagnetic stress will not occur in isolation. Instead, biological systems will experience a unique triple-stressor environment consisting of:
\begin{itemize}
    \item Hypomagnetic field deprivation ($< \SI{20}{\nano\tesla}$ across northern candidate sites)
    \item Reduced partial gravity ($0.38\,g$)
    \item Chronic galactic cosmic radiation (GCR) and solar particle events (SPE), yielding surface doses of $\sim 200\text{ to }\SI{300}{\milli\sievert}/\text{year}$
\end{itemize}

Crucially, experimental evidence suggests potent synergy among these stressors. In rodent models, hypomagnetic field exposure dramatically aggravates musculoskeletal bone loss induced by mechanical unloading \citep{Xu2012, Mo2014}. Furthermore, HMF-induced perturbations in radical pair spin-state recombination kinetics and ROS accumulation \citep{Binhi2017, Hore2016} can impair cellular DNA repair pathways, potentially sensitizing biological systems to ionizing radiation damage.

\subsection{Operational Recommendations for Mars Surface Payloads}
To de-risk future human exploration and ensure the viability of long-duration surface greenhouses, we propose three core recommendations:
\begin{enumerate}
    \item \textbf{Surface Magnetometry Integration:} All future landers and rover missions must incorporate calibrated fluxgate magnetometers to expand ground-level magnetic mapping beyond the single InSight measurement.
    \item \textbf{Active Magnetic Modulation in Growth Chambers:} Agricultural and bioregenerative life support modules deployed in northern lowlands should integrate Helmholtz coil arrays to generate artificial Earth-like static magnetic fields ($\sim\SI{50}{\micro\tesla}$), restoring normative cryptochrome signaling, auxin polar transport, and optimal crop yields.
    \item \textbf{Controlled Magnetic-Biology Payloads:} Early robotic precursor missions should conduct dual-chamber biological experiments (shielded near-null field vs. \SI{50}{\micro\tesla} artificial field) directly on the Martian surface to empirically isolate the biological consequences of the Martian magnetic environment.
\end{enumerate}
